\label{chap:discussion}

The segmentation numbers in Chapter~\ref{chap:results} are strong within the repeated Bast\o{} route, but the weakest cases are not a single failure mode. The Jan 2024 development split holds short runs of large model-label disagreement, the Oct 2024 validation split concentrates its residual disagreement in shallow near-stationary recordings, and the repeated-crossing analysis finds its largest cross-pass spread in the deep channel and adjacent slope. The first question is what the validation can establish without an independent seabed-depth reference.

\section{Evaluating without independent reference labels}
\label{sec:no-reference-labels}

The labels carry no independent measurement of seabed depth. They come from the algorithmic label pipeline and are used as it produces them, so the training target is an estimate rather than a measured truth. In diagnostic-test terms, the study lacks a gold standard \cite{rutjes_evaluation_2007}. The metrics in Chapter~\ref{chap:results} therefore measure agreement with the label pipeline, not distance from the true seabed: an IoU of 0.98 says the model reproduces the reference, not that the reference was correct.

This would matter less if the labels were wrong in a random, uncorrelated way, which risk minimisation can tolerate for suitable losses \cite{manwani_noise_2013}; label-noise reviews separate that simplest regime from class-dependent and feature-dependent noise \cite{karimi_deep_2020}. The errors here are the harder kind: phantom returns from single-beam sidelobe artefacts, second-echo confusion at twice the true depth, and bubble reverberation in the harbour approaches. A model trained to match those labels is rewarded for reproducing the same systematic mistakes, and the bias does not wash out with more data. Two consequences follow: a high agreement score cannot separate a correct model from one that has copied a biased reference, and model-label disagreement is not automatically a model error.

The same limitation applies to supervised conformal coverage. The nonconformity score is computed from the model-reference depth discrepancy, so the interval is calibrated to contain the algorithmic reference bottom, not the physical seabed. A band can miss because the model rejects a bad detector pick, and it can cover because model and reference share the same bias, so reference coverage is an internal diagnostic rather than a measure of operational accuracy.

The per-file results separate two observations. First, the 57 trained models rank Oct 2024 validation files by IoU almost identically, with a mean pairwise Spearman correlation of 0.895. This agreement indicates that difficulty is concentrated in the same files across model variants, but it does not identify whether the disagreement originates in the models, the reference labels, or both. Second, visual inspection of selected difficult files shows cases where the model line follows the visible seabed return more closely than the reference (Figures~\ref{fig:grid-worst-max} and~\ref{fig:grid-worst-median}; Appendix~\ref{app:prediction-grids}). These examples support reference-label error as one explanation for some low scores, without establishing its frequency or excluding model error. Numerical penalty and visual plausibility can therefore point in opposite directions, but the visual assessment is selective and subjective.

Separating agreement from accuracy requires a reference partly independent of the detector: blinded expert picks, co-located bathymetry where it exists, or a latent-class treatment over several detectors \cite{hui_estimating_1980,rutjes_evaluation_2007}. None is available here, and the headline numbers should be read against that limit.

\section{Primary within-route performance}

The headline IoU is high partly because it scores the filled bottom-or-below mask. After thresholding at 0.5, a one-pixel boundary error is negligible in an area metric over a deep echogram, so the model is graded weakly on boundary sharpness. The line-mask variant reaches comparable IoU but worse depth deviation, because a small misalignment in a thin line propagates directly into the depth metric while the area mask absorbs it (Section~\ref{sec:dataset-results}). The residual model-reference difference is diffuse: the median per-file depth deviation is near 0.3\,m across most files (Section~\ref{sec:segmentation-performance}). Pixel discretisation, model error, and variation in the algorithmic reference could all contribute to this difference; the present experiments do not separate them. The mean IoU of 0.9972 is therefore a measure of overlap with the detector-derived reference, not a lower bound on error or a distance from the physical seabed.

The per-file results show a ceiling with a short tail (Table~\ref{tab:per-file}): eleven files exceed 0.99 IoU, 13 have sub-metre median depth deviation, one has a narrow 5\,m-scale offset, and five fall to 0.927--0.971 IoU with mean deviations of 6.4--18.8\,m. The low-IoU files may be label-noise cases, as Section~\ref{sec:no-reference-labels} suggests, or conditions the model has not learnt. Their within-file structure favours localised failure: D20240131-\allowbreak T211740 has a median deviation of 0.95\,m but a 95th percentile of 92\,m. Checking whether those files coincide with near-dock, low-speed segments (Section~\ref{sec:vessel-speed-prior}) would separate the two readings, but was not run.

\section{Model selection and diagnostics}

Region IoU was insensitive to several design choices within the Jan 2024 development subset: five of the six dataset variants fall within 0.4~pp of the baseline, the loss variants span 1.05~pp, and the three ResNet-34 freezing strategies span 0.34~pp. Easy interior pixels dominate the area metric while the relevant variation sits in a thin boundary band, so a boundary-focused metric would retain differences that pooled region IoU absorbs \cite{cheng_boundary_2021}. A separate augmentation ablation used a 100-file random subset, with 80 files for training and 20 for validation (Section~\ref{sec:augmentation-results}). Within that protocol, IoU stayed above 0.99 across all seven configurations, while structured corruptions reduced depth deviation from 2.49\,m to 0.80\,m and tightened the smallest across-file IoU spread to 0.0025. The gains were associated with channel dropout and ping-stripe replacement rather than flips and photometric shifts. These corruptions were designed to mimic dropped frequencies and lost pings, thereby encouraging the model to use neighbouring spatial context instead of relying on one intact column. The ablation results are consistent with that mechanism, but they do not show which features the augmented models used. Since the augmentation and Jan 2024 comparisons use different subsets, their effect sizes cannot be ranked against one another. Training-set size was not varied, so the augmentation experiment also does not establish whether the benefit depends on data scarcity. It only shows that structured augmentation improved depth stability on its fixed 80-file training subset.

The convergence curves require the same protocol distinction. Training metrics are averaged over batches after channel dropout and ping-stripe corruption, whereas validation metrics are computed on unaugmented patches. Higher validation IoU and lower validation loss can therefore arise without superior validation performance on comparable inputs. Running averages during early epochs may widen the visual separation because they combine predictions from several weight states. The curves remain close through the selected epoch and show no sustained divergence associated with memorisation.

Replacing ResNet-34 with ResNet-50 did not improve accuracy (Section~\ref{sec:transfer-learning-results}), but this says little about encoder depth given the small training set used for that experiment. The relevant observation is that ImageNet initialisation stayed competitive on acoustic power images, consistent with generic early convolutional features \cite{yosinski_how_2014}. The spread across freezing strategies is narrow, 0.9781 for the frozen encoder against 0.9812 for full fine-tuning, and the Bast\o{} Standard U-Net (\texttt{chunk\_arch\_sweep\_2}), trained without ImageNet pretraining, still reached 0.9972 IoU on the held-out Oct 2024 period. Either min-max normalisation to $[0,1]$ wasted part of the ImageNet statistics, or the task needs little of ImageNet's higher-level invariance; a direct check would cross random against ImageNet initialisation with min-max against ImageNet-style standardisation under one split and freezing schedule.

\section{Uncertainty evaluation}

The uncertainty results give a specific, limited advantage over the traditional detector. On held-out Oct 2024 files the probability map is sharply bimodal, with ambiguous pixels concentrated in a sub-metre boundary band rather than spread through the water column. Its raw calibration is weaker than the expected calibration error of $0.0007$ suggests: the reliability curve is mostly underconfident, then overconfident only in the highest bins. This local pattern differs from the usual overconfidence reported for accurate deep networks with sigmoid or softmax outputs \cite{guo_calibration_2017}, but it leads to the same operational caution. The 50/50 BCE-Dice loss, IoU checkpointing, class imbalance, and filled-mask representation are possible contributors because they prioritise segmentation overlap rather than probability calibration. Their effects were not isolated experimentally, and the concentration of pixels in the confident tails also makes the aggregate ECE insensitive to the narrow boundary region. The sigmoid output should not be treated as the main uncertainty guarantee.

Split-conformal calibration is more useful, once the target of the guarantee is stated precisely. A constant-width band fitted on the scored Jan 2025 conformal cache covers the corresponding scored Feb/\allowbreak Jun 2025 cache at $92.6\%$ for a nominal $90\%$ target, and the supervised locally weighted band is similarly conservative at $93.9\%$. These numbers say the model-reference discrepancy is conservatively calibrated on this within-route cache, not that the interval contains the true seabed with the same probability. Marginal coverage also does not show that the probability-transition width identifies difficult individual pings. The scale uses the first $p=0.25$ and $p=0.75$ crossings in the full water column, so unrelated probability structure above the predicted bottom can affect it. The constant band therefore remains the clearer supervised uncertainty budget.

The ensemble results fit this. Same-architecture seed averaging changed neither the point prediction nor the adaptive band beyond the small differences in Section~\ref{sec:ensemble-results}, because the members made nearly the same errors rather than cancelling one another. A separate bottom-anchored transition width correlated more strongly with error and produced a narrower mean band on the full February-June diagnostic split, but that experiment changed the point predictor, transition definition, scale floor, and evaluated files. It cannot be read as an improved version of the scored-cache interval without a common comparison. The correlated errors indicate a shared component under the present architecture, training data, and labels, but they do not show that this component is irreducible. Possible sources include common reference-label bias, model bias, and acoustic ambiguity. Ensemble disagreement is too flat to act as a local error scale in these experiments, so the learned detector provides boundary uncertainty and conformal depth intervals for the reference-label problem rather than a universal failure alarm.

Patch-scale inspection of the conformal examples (Section~\ref{sec:patch-conformal-examples}) shows the adaptive band tracking local difficulty directly, which the pooled coverage number cannot. On a near-flat shallow bottom the locally weighted band can contract below the residual jitter of an easy boundary, missing many reference pings while the bottom-line error stays near a single sample, so low depth deviation and a high patch-scale miss count are not in conflict (Figure~\ref{fig:patch-conformal-flat}). Where the model and reference part company, the same band widens over the central disagreement and holds the reference inside, recovering coverage that the fixed-width band loses (Figure~\ref{fig:patch-conformal-hard}). Band width then carries information the point estimate does not, and in the useful direction: it opens where the local estimate is least trustworthy. The adaptive band is in this sense a diagnostic rather than a guarantee. Its width can be read off the echogram during inspection, or thresholded as an automatic flag for segments that warrant a second look. The easy-patch misses and the moderate per-ping correlation between width and error confirm that it is not a calibrated per-ping interval.

The adaptive result that earns emphasis is the label-free repeated-crossing band. Its local scale is measured repeatability of the same route bin across many crossings, not a weak model proxy and not the traditional detector's per-ping agreement score. It preserves the $90\%$ target while reducing median half-width from $7.59$\,m to $3.59$\,m and widening mainly in empirically inconsistent deep-channel bins, identifying where a repeated route needs a wider uncertainty budget without manually labelled seabed depths. The GPS positions carry their own noise (Chapter~\ref{chap:data}), possibly larger on the deeper crossings, so part of the deep-channel spread read here as seabed inconsistency may instead reflect uncertainty in where each crossing was located.

\section{External transfer}
\label{sec:generalise-beyond}
The model was trained and evaluated only on the Bast\o{} Electric route. A single fixed installation made the pipeline coherent but left geographic transferability largely untested, and integrating other campaigns would require harmonising transducer configuration, pulse settings, calibration state, and file conventions. The zero-shot hake diagnostic (Section~\ref{sec:transfer-survey}) shows the limit directly: without fine-tuning the model places the seabed far from the algorithmic line along the slope and marks dense pelagic scattering layers as bottom where the seabed lies beyond the 200\,kHz range, and agreement against the 38\,kHz detector stays moderate at 0.61--0.84 IoU. These scores measure agreement with an algorithmic reference rather than verified accuracy, and the three transfer experiments below use disjoint file splits so their results do not overlap.

Threshold recalibration does not improve the zero-shot point prediction. It restores marginal coverage against the hake detector by increasing the fitted quantile, and the cost is the constant-width median full width rising from 4.2\,m under Bast\o{} calibration to 42.3\,m under hake calibration (Table~\ref{tab:hake-recalibration}). The low zero-shot coverage therefore reflects a mismatch between the Bast\o{} calibration scores and the hake model-reference discrepancies, while the recalibrated result trades interval width for reference coverage. Neither result establishes coverage of the physical seabed. Calibration data must consequently represent the intended survey and model, and the resulting width must remain operationally acceptable.

Fine-tuning the Bast\o{} Standard checkpoint on a hake training split (Table~\ref{tab:hake-finetune}) raised held-out hake agreement substantially and more than halved mean depth deviation, and the Bast\o{}-calibrated constant-width band recovered from 7.9\% to 89.8\% coverage on the adaptation groups. Training the same architecture from scratch on the hake split went further still, which shows the hake detector labels are learnable from the target survey itself. That gain does not point to one universal checkpoint. On the Bast\o{} test split the fine-tuned model held its accuracy at 0.989 IoU against the 0.990 of the Bast\o{} Standard U-Net, but the hake-only scratch model collapsed to 0.171 IoU and a 133.5\,m mean depth deviation. The practical reading is domain-specific deployment. If the operating survey is Bast\o{}, the Bast\o{} Standard checkpoint remains the right model; if it is hake-like, a small labelled or detector-labelled in-situ set can adapt the model and recalibrate the uncertainty bands, though the adapted checkpoint must be validated separately before use elsewhere.

\section{Limitations}
\label{sec:weaknesses}

The largest practical weakness is the same one that limits interpretation in Section~\ref{sec:no-reference-labels}: label quality. The detector fails in regimes where it cannot reliably flag its own error, so automatic cleaning is unreliable. A follow-on study should keep the pipeline only if it also adds an independent reference, expert review, or an exclusion rule for ambiguous segments.

The qualitative examples suggest the model sometimes resists label errors by placing the seabed where the echogram supports it, but they do not measure how often this happens. Phantom returns, second-echo picks, and residual deep-water disagreement were observed during selective inspection, but their prevalence among the roughly 11{,}000 patches is unknown. The data therefore do not establish whether gross errors are rare outliers or whether any residual label noise is close to the benign class-independent regime described by \citet{manwani_noise_2013} and \citet{karimi_deep_2020}. A per-condition label-error estimate is needed to determine which mechanism dominates. The channel ablation (Section~\ref{sec:channel-importance}) sharpens the concern: the label pipeline runs on 38\,kHz alone, and that same channel carries most of the network's predictive signal. A high score may therefore reflect reproduction of the 38\,kHz detector rather than an independent seabed representation. Separating the two needs a target the 38\,kHz detector did not produce, such as the expert-labelled reference proposed in Section~\ref{sec:expert-labelled-reference}.

\subsection{Development ablations on a single recording period}
The model-comparison ablations in Tables~\ref{tab:transfer}, \ref{tab:loss}, and \ref{tab:datasets} use 89 files from 31 January to 1 February 2024, split within that recording period. This prevents patch leakage but keeps weather, seasonal biology, instrument state, and route conditions nearly fixed, so the tables are controlled development comparisons, not thesis-wide generalisation estimates. Later Oct 2024 validation, Jan 2025 calibration, Feb/\allowbreak Jun 2025 tests, and repeated-crossing analyses mitigate this for the final model but do not turn the Jan 2024 ablation gaps into across-season effects.

\subsection{Mask labelling}
The training masks were used as the 38\,kHz label pipeline produced them, without systematic correction or filtering of visibly wrong patches. Multi-frequency averaging or detector ensembling could disambiguate some cases, but this proof of concept kept label generation to the single channel that carries the seabed return in deep water. Visual inspection was necessarily selective: the working set holds roughly 181\,GB of raw files and over six million pings, yielding about 11{,}000 patches, so exhaustive review was infeasible. Spot checks of Bast\o{} outputs and hake transfer examples found traditional-detector failures and phantom-return labels, leaving an unknown mislabelled fraction in both training and test sets. Some failures leave no reliable visual trace, because single-beam sidelobe aliasing and diffraction can produce returns an operator cannot separate from the true seabed on the echogram alone.

The repeated-transect design gives a quantitative handle on label noise. A fixed seabed point is recorded on hundreds of crossings, so assigned depth should be stable; where it is not, the label is unreliable even if any single echogram looks plausible. Section~\ref{sec:label-consistency} shows robust standard deviations near 50\,m at the least consistent bins, with individual crossings disagreeing by more than 100\,m and some bins splitting into two clusters. These bins lie mainly in the deep central channel and adjacent slope, so instability is not confined to harbour approaches; the repeated crossings map label failure better than they supply a clean depth reference.

The harbour approaches at Moss and Horten contain some of the clearest visually identified detector failures, where bubble-like layers occur near a shallow seabed and the traditional detector selects them as bottom. Braking or bow-thruster activity is a plausible source of these layers, but vessel operations were not recorded as explanatory variables and no causal association was tested. The interpretation therefore remains a visual hypothesis. Visual filtering catches egregious failures but cannot recover a stable reference where the signal is ambiguous; cleaning this regime requires a pipeline that detects disagreement between detectors, or between detector and echogram, and repairs or excludes affected segments. At the other extreme, some archive segments were recorded while the ferry was moored overnight. These could be excluded by service-hour filtering, but inspection showed them among the most reliable labels, with a stable seabed return the detector tracks easily, so we kept them as low-noise examples. No ablation removes docked recordings, so their effect on the final model remains unmeasured. Appendix~\ref{app:time-of-day} reports hourly coverage and midnight examples.

\subsection{Correlation and exchangeability in the conformal calibration}
\label{sec:conformal-exchangeability}
The supervised conformal statement in Section~\ref{sec:conformal-theory} is a marginal guarantee for an exchangeable calibration-test sequence, and the implemented score treats each ping as one example. Two features of the data strain that interpretation. Consecutive pings are not independent draws: adjacent transmissions illuminate nearly the same seabed, so nearby nonconformity scores move together and the informative sample size is closer to the number of calibration files than the number of pings. The file-level split compounds this, assigning entire days, sea-states, and tidal phases to one side of the split; the repeated-transect design makes block exchangeability defensible within the route, not exact, and a systematic shift in depth-deviation distributions biases $\hat q$ in a direction the marginal guarantee cannot register \cite{tibshirani_conformal_2019}. The correction is to treat the file, not the ping, as the diagnostic unit. Per-file coverage, or a conformal variant for dependent or non-exchangeable observations \cite{chernozhukov_exact_2018,barber_conformal_2023}, would expose the effective sample size a ping-level count hides.

The diagnostics in Section~\ref{sec:conformal-exchangeability-results} bear this out only in part. Calibration and test score distributions in the scored conformal cache differ by a Kolmogorov-Smirnov distance near $0.09$, far above the within-calibration baseline of about $0.005$, yet realised coverage is conservative at $92.6\%$ against the $90\%$ target. This is narrow reassurance, because one conservative draw is still compatible with high variation across recalibrations. The file-level breakdown makes the effective sample size visible: the pooled $92.6\%$ separates into $95.6\%$ on one test period and $89.8\%$ on the other, while the unweighted file mean is $93.8\%$ with standard deviation $0.074$. Leaving one calibration file out moves $\hat q$ with a coefficient of variation of $10.2\%$, so the threshold behaves more like it was set by 29 files than by $80{,}180$ pings. Reporting realised coverage with file-level diagnostics better reflects the calibration evidence than a pooled ping-level number alone.

\subsection{Asymmetric area of interest when running live}
Training and offline evaluation use patches where the ping of interest sits inside the window, with context on both sides. Live deployment is different: the newest ping is the right-most column, with no future context, so a model trained on symmetric patches sees a different effective context at the live edge. The per-column analysis in Appendix~\ref{app:edge-degradation} shows a real but bounded cost. Over the rightmost 5\% of a patch IoU drops by only 0.0049 from the interior, but the final column falls from 0.9829 to 0.9596 and bottom-line error on the resized grid more than doubles, from 13.5 to 32.4 rows. A live detector reporting only the newest column would operate at the weakest part of the patch, and the worse left-edge behaviour suggests padding, reduced context, and resizing artefacts also contribute. Mirroring around the current ping would restore the input statistics but adds no future information and may invent continuity at abrupt depth changes; the direct options are to train in the causal geometry, with the target column at the live edge, or to delay reporting until the ping has moved into the patch interior when latency allows.

\subsection{Vessel speed, motion, and the spatial prior}
\label{sec:vessel-speed-prior}
The convolutional model treats the ping axis as a regular spatial grid. If vessel speed varies while the ping interval remains fixed, along-track spacing also varies and the same physical seabed length occupies a different number of columns. Such compression is plausible during harbour approaches, but vessel speed was not related to prediction error in the present analysis, so its contribution remains a hypothesis. Vessel motion introduces a second possible installation dependence. The pipeline treats sample index as range below the transducer without explicit motion compensation, and heave, pitch, or roll could alter the apparent seabed trace across pings. These effects were not measured here. The reported performance is therefore conditional on the Bast\o{} Electric installation and route; transfer to vessels with different motion characteristics would require motion data and a direct test of the relevant corrections.

\subsection{Dependence on dual-frequency input}
The patch representation stacks 38\,kHz, 200\,kHz, and their difference, so inference expects two co-registered frequencies and a single-frequency echosounder or different pair would require retraining. As Section~\ref{sec:channel-importance} shows, 38\,kHz carries most of the signal: a 38\,kHz-only model reaches 0.99 IoU, and removing 38\,kHz is the most damaging change. The difference channel contributes little. It is formed after the two frequencies have been normalised separately, so it omits the absolute dB offset between bands and is a linear combination of inputs already supplied to the network. This construction offers one explanation for the ablation result, but the experiment does not isolate the effect of normalisation from redundancy between channels. A dedicated single-frequency 38\,kHz model is therefore a plausible simplification.

\subsection{Single-bottom assumption}
The mask labels the bottom-or-below region as one class and assumes each ping carries exactly one monotonic boundary. Sub-bottom layering, fish resting on the seabed, and second-bottom echoes from multipath are not represented and cannot be recovered as distinct structures. The crossing rarely violates this, but the formulation breaks on terrain with overhangs, gas seeps, or strong volume scattering near the floor.

\subsection{Absent-bottom evaluation}
Fully empty reference files are excluded from the supervised headline metrics, so the thesis does not establish how reliably the model identifies that no seabed is present within the recorded range. This matters most under transfer: some hake recordings contain no seabed return, yet the model can still assign bottom probability to scattering layers or open water. These cases are a deployment limitation rather than part of the reported segmentation averages. A future evaluation requires an explicit empty-target rule and a separately reported absent-bottom detection metric.

\subsection{Patch resizing and range resolution}
The 89-zarr baseline used 256-ping patches resized to $512\times512$, distorting the physical aspect ratio. The main model extracts $512 \times 512$ patches without along-track interpolation, resampling range only. The square target is a modelling choice, not a ResNet-34 or ImageNet constraint, since the encoder can process other sizes as long as dimensions stay divisible by 32. The hake transfer estimates the cost: stretching a 214-ping file to the 512-ping window drops mean IoU by about 0.06 against stitched native-width inference (Section~\ref{sec:transfer-survey}), so native-resolution training and aspect-ratio-preserving targets remain untested.

Range resampling also lowers vertical resolution whenever a patch holds more range samples than the 512-row target. The number of range samples and the maximum recorded range can differ across files, so one output row has no common conversion to metres throughout the dataset. The present experiments do not isolate how much this resampling contributes to bottom-depth error or whether its effect differs between shallow and deep recordings. The two-stage refinement of Section~\ref{sec:highres-refinement} would test whether higher local range resolution improves boundary localisation.

\section{Future Work}

\subsection{Expert-labelled training and test data}
\label{sec:expert-labelled-reference}
The next detector should start from a better reference. An expert-picked subset, ideally blinded to the algorithmic output, would supply less noisy training examples and a test set where agreement begins to approximate accuracy rather than reproduction of the labeller. Phase-derived split-beam angle features extracted from the raw IQ recordings are a candidate auxiliary input for filtering or down-weighting ambiguous labels, but their value should be tested against expert picks because distributed seabed echoes do not yield a clean single-scatterer angle (Section~\ref{sec:beam_formation}). The harbour approaches at Moss and Horten should be prioritised because bubble layers and shallow bottoms defeat the traditional detectors.

\subsection{Stress testing under adverse conditions}
The validation set includes steep slopes and harbour reverberation, but a systematic stress test should deliberately sample abrupt depth transitions, heavy bubble sweep-down, and strong biological scattering near the bottom. Curated difficult segments would show where the model breaks and whether uncertainty widens appropriately, and a linear mixed-effects model with covariates such as depth, local slope, and signal-to-noise ratio could then identify the factors most predictive of failure.

\subsection{Generalisation to other surveys}
Training across vessels, transducer configurations, and regions is the next generalisation test. Given the harmonisation effort described in Section~\ref{sec:generalise-beyond}, fine-tuning the Bast\o{} Electric model on smaller labelled sets from other surveys is likely more practical than training from scratch on a merged dataset.

\subsection{Temporal modelling across pings}
The current model processes patches independently and recovers along-track continuity only through convolutional context inside each window. A recurrent or attention layer over the ping axis \cite{shi_convolutional_2015,vaswani_attention_2017}, or a temporal smoothing penalty, would encode sequential context directly. The harbour approaches are the most likely beneficiaries, but the gain must be measured.

\subsection{Denser patch sampling}
Overlapping training patches would multiply the patch count, but mostly by adding correlated views of the same seabed. The roughly 11{,}000 existing patches already cover the route densely, and compute rather than patch count was the binding constraint. The training and validation loss curves remain close through the selected epoch (Section~\ref{sec:convergence}), so they provide no evidence that denser sampling is needed to close an optimisation gap. Whether a finer stride improves boundary localisation or instead adds redundant examples remains untested.

\subsection{Domain-adapted encoder pretraining}
Self-supervised pretraining on unlabelled EK80 recordings could provide an encoder closer to echograms than ImageNet and lessen dependence on algorithmic labels, following domain-specific pretraining in other imaging domains \cite{chen_simple_2020,azizi_big_2021}. This pairs naturally with transfer to other surveys.

\subsection{Uncertainty validity under distribution shift}
The conformal bands rely on calibration-test exchangeability, which is plausible within the repeated route but not across vessels or regions (Section~\ref{sec:conformal-exchangeability}). In the hake transfer, Bast\o{}-fitted bands fall far below the $90\%$ target for the constant-width and locally weighted bands: zero-shot held-out coverage is 20.4\% and 30.5\%, respectively (Section~\ref{sec:transfer-survey}). A small in-situ calibration set repairs the marginal constant-width coverage estimate, which suggests practical domain adaptation is possible, but it also exposes the cost: the constant-width band expands to a 42.3\,m median width. The fine-tuning experiment adds a complementary route: changing the model weights first reduces the residual shift, after which Bast\o{} or balanced pooled calibration can again give near-nominal hake coverage with narrower intervals. Whether this remains true for a broader set of vessels, surveys, and independently checked seabed labels remains open.

\subsection{High-resolution bottom refinement}\label{sec:highres-refinement}
A two-stage approach could improve boundary accuracy: first run the current model at full scale, then crop a narrow vertical window around that estimate and refine it at higher resolution. This concentrates capacity at the seabed without processing the full depth range at high resolution.
